\chapter{Краткое описание алгоритма коллаборативной фильтрации}

Алгоритм коллаборативной фильтрации прогнозирует оценку целевого пользователя для объекта на основе оценок схожих пользователей. Сходство между пользователями вычисляется по их векторам оценок с использованием мер близости, таких как косинусное сходство или корреляция Пирсона. Для целевого пользователя идентифицируется $k$ ближайших соседей --- пользователей с наивысшим показателем сходства. Итоговая прогнозируемая оценка рассчитывается как взвешенная средняя оценок этих соседей, где весами выступают коэффициенты их сходства с целевым пользователем. Рекомендации формируются путем ранжирования объектов с наибольшими прогнозируемыми оценками.

\clearpage
